\chapter{Demo}
\section{A section}

This is the first chapter of this thesis' main part. Page numbering switches over from roman to arabic. 

\minisec{Citations}
Citations can be done like this: \cite{Amann2008AnalysisII}. The prefered citation style may be set in \texttt{Options.tex}.

\minisec{Quotes}
Quotes can be typeset as follows: \enquote{This is a quote}. They will be automatically typeset according to the chosen language.

\minisec{Enumerations}
\begin{enumerate}[wide]
	\item For enumerations that need not be particularly exposed (e.g. in an assertion type statement as inside a lemma, a proposition, a theorem or a corollary or in a definition), the option \enquote{wide} as set for
	this enumeration is preferable. It allows for better utilisation of the type are.
	\item For enumerations as listed above, the option \enquote{wide} is probably not best suited.
\end{enumerate}

\minisec{Special environments and references}
Special environments can be typeset as follows:

% Aside of Theorem you can also use Proof, Lemma, Definition, Example, Corollary, Proposition, Construction, equation, IEEEeqnarray, Remark, Reminder.
\begin{Theorem}[Heine-Borel]
\label{thm:heine_borel}
	A subset $A\subseteq \mathds{R}^n$ is compact if and only if it is closed and bounded. 
\end{Theorem}
You can reference them as follows: \cref{thm:heine_borel}. If you just want the number and not the type of the referenced object, you can also use \ref{thm:heine_borel}.